\documentclass[a4paper]{article}     
\usepackage{amsfonts}
\usepackage{amsmath}
\usepackage{amssymb}
\usepackage{graphicx}
\usepackage{color}

\newcommand{\Q}{\mathbb{Q}}
\newcommand{\R}{\mathbb{R}}
\newcommand{\llim}{\lim\limits}
\newcommand{\dint}{\displaystyle\int}

\begin{document}

\noindent \large Auteur: Hina Rana \\
\noindent \large Studierichting: Informatica\\
\noindent \large Oefeningenreeks 4A nr. 37.9\\

\medskip

\normalsize

$ \dint \dfrac{7}{x^2 + 8x + 18} dx $ \\

$ = 7 \dint \dfrac{1}{x^2 + 8x + 18} dx $ \\

$ = 7 \dint \dfrac{1}{(x+4)^2 + 2} dt $ \\

stel $t = x+4 \Rightarrow dt = dx$ \\

$ = 7 \dint \dfrac{1}{t^2 + 2} dt $ \\

$ = 7 \dint \dfrac{1}{2 \left(\dfrac{t^2}{2} + 1 \right)} dt $ \\

$ = \dfrac{7}{2} \dint \dfrac{1}{\left(\dfrac{t}{\sqrt{2}} \right)^2 + 1} dt $ \\

stel $u = \dfrac{t}{\sqrt{2}} \Rightarrow dt = \sqrt{2} du$ \\

$ = \dfrac{7}{\sqrt{2}} \dint \dfrac{1}{u^2 + 1} du $ \\

$ = \dfrac{7}{\sqrt{2}} \operatorname{Bgtan} \left(\dfrac{t}{\sqrt{2}} \right) + C $ \\

$ = \dfrac{7}{\sqrt{2}} \operatorname{Bgtan} \left(\dfrac{x+4}{\sqrt{2}} \right) + C $\\


\begin{thebibliography}{999}
\end{thebibliography}
\end{document} 
